\chapter{Apêndice ou Anexo?}
\label{ap:diferenca}

O Apêndice~\ref{ap:diferenca} explica a diferença entre apêndice e
anexo, conforme as convenções do \emph{template} oficial da
\thesisuniversity.

\paragraph{Apêndice.} Apêndices englobam materiais elaborados pelo
autor(a), tais como gráficos, quadros, tabelas, traduções,
organogramas e esquemas que prestem informação relevante para a
compreensão do trabalho. Só devem figurar nos apêndices as
informações previamente referenciadas no texto. As informações são
total ou parcialmente da responsabilidade do autor.

\paragraph{Anexo.} Anexos englobam documentos que, não sendo
elaborados pelo autor, serviram de base para a construção do estudo
ou facilitam a compreensão da tese/dissertação. Só devem figurar
nos anexos documentos e/ou materiais previamente referenciados no
corpo do trabalho. Todos os documentos devem estar em formato
digital.

% ============================================================================
\chapter{Cálculo de Spearman e Kendall em Python puro}
\label{ap:correlation-code}

Este apêndice apresenta a implementação completa do \emph{script}
\code{compute\_rank\_correlations.py} (Capítulo~\ref{ch:metodologia},
Secção~\ref{sec:metricas}), em Python puro e sem dependência de
\code{scipy}. O \emph{script} consome os ficheiros
\code{ranking\_results.json} produzidos pelo \code{ranker.py} e
gera o ficheiro consolidado \code{rank\_correlations.json},
utilizado para popular a Tabela~\ref{tab:correlacoes}.

\begin{lstlisting}[language=Python, caption={Implementação de Spearman $\rho$ e Kendall $\tau_b$ em Python puro (excerto).},label={lst:rank-code},basicstyle=\ttfamily\scriptsize]
def mean_ranks(scores):
    """Converte scores em ranks com mean-rank para ties."""
    indexed = sorted(enumerate(scores), key=lambda x: -x[1])
    ranks = [0.0] * len(scores)
    i = 0
    while i < len(indexed):
        j = i
        while j + 1 < len(indexed) and indexed[j+1][1] == indexed[i][1]:
            j += 1
        avg = (i + 1 + j + 1) / 2.0
        for k in range(i, j + 1):
            ranks[indexed[k][0]] = avg
        i = j + 1
    return ranks

def spearman_rho(a, b):
    return pearson(mean_ranks(a), mean_ranks(b))

def kendall_tau_b(a, b):
    """Kendall tau-b com correcao para empates (Goodman, 1954)."""
    n = len(a); C = D = Ta = Tb = 0
    for i in range(n):
        for j in range(i + 1, n):
            da, db = a[i] - a[j], b[i] - b[j]
            if   da == 0 and db == 0: Ta += 1; Tb += 1
            elif da == 0:             Ta += 1
            elif db == 0:             Tb += 1
            elif (da > 0) == (db > 0): C += 1
            else:                      D += 1
    P0 = n * (n - 1) / 2
    den = ((P0 - Ta) * (P0 - Tb)) ** 0.5
    return (C - D) / den if den else float('nan')
\end{lstlisting}

A versão completa, incluindo o carregamento dos ficheiros JSON e a
escrita do resultado consolidado, encontra-se em
\code{compute\_rank\_correlations.py} na raiz do repositório
público.

% ============================================================================
\chapter{Mapeamentos \code{func\_to\_req} dos três casos de estudo}
\label{ap:mapeamentos}

Os mapeamentos \code{func\_to\_req} utilizados em cada caso de
estudo encontram-se em:

\begin{itemize}
    \item CPython \emph{stdlib}: \code{testes/simples\_stdlib/mapeamento.py}
    \item Flask: \code{testes/medio\_flask/mapeamento.py}
    \item \emph{scikit-learn}: \code{testes/complexo\_sklearn/mapeamento.py}
\end{itemize}

Os ficheiros são distribuídos como parte do repositório público
mencionado em Capítulo~\ref{ch:introducao},
Secção~\ref{sec:contribuicoes}. Cada ficheiro contém um único
dicionário Python \code{FUNC\_TO\_REQ} cujas chaves são nomes
qualificados de funções/métodos e cujos valores são identificadores
de requisitos no formato \texttt{REQ\_*}.
